\section{The stellar radiation field and its absorption}
\label{sec:radiation}

\subsection{The incident spectrum}
\label{sec:radiation:spectrum}

The star is represented by a plane beam that falls radially onto the
substellar column. Its spectral shape is chosen once, by the
\texttt{Spectrum type:} key of \texttt{input.inp}, and that single choice
fills \emph{every} band of the photon grid: the extreme ultraviolet and
X-ray bands, and also the part below the hydrogen Lyman edge where the
He\,$2^3$S metastable, the low-ionization-potential metals and hydrogen in
$n=2$ absorb. The three types are implemented in
\texttt{J\_inc.f90} (\texttt{stellar\_flux\_eV}) and
\texttt{sed\_read.f90}:

\begin{description}
\item[Power law] (\texttt{is\_PL\_sed}). Two power laws of one common index
  $\alpha \equiv$ \texttt{PLind}, normalized separately on the EUV interval
  $[E_{\rm low},E_{\rm mid}]$ and on the X-ray interval
  $[E_{\rm mid},E_{\rm top}]$. With
  $L_{\rm X}$ and $L_{\rm EUV}$ the two band luminosities, $a$ the orbital
  distance and $\mathcal{R} = 10^{L_{\rm X}-L_{\rm EUV}}$,
  %
  \begin{equation}
  J_{\rm XUV} = \frac{10^{L_{\rm X}} + 10^{L_{\rm EUV}}}{4\pi a^{2}},
  \qquad
  J_{\rm X} = \frac{\mathcal{R}}{1+\mathcal{R}}J_{\rm XUV},
  \qquad
  J_{\rm EUV} = \frac{1}{1+\mathcal{R}}J_{\rm XUV},
  \label{eq:rad:bandsplit}
  \end{equation}
  %
  and, for $\alpha \neq -1$ and $P_1 = \alpha+1$,
  %
  \begin{equation}
  J_{\rm inc}(E) =
  \begin{cases}
    \dfrac{P_1 J_{\rm EUV}}{E_{\rm mid}^{P_1}-E_{\rm low}^{P_1}}\,E^{\alpha},
      & E < E_{\rm mid},\\[2ex]
    \dfrac{P_1 J_{\rm X}}{E_{\rm top}^{P_1}-E_{\rm mid}^{P_1}}\,E^{\alpha},
      & E \ge E_{\rm mid},
  \end{cases}
  \label{eq:rad:powerlaw}
  \end{equation}
  %
  the logarithmic form replacing the ratio at $\alpha=-1$, and the X-ray
  normalization set to zero when \texttt{thereis\_Xray} is false. Below
  $E_{\rm low}$ the same law is extrapolated, which is what the type means
  rather than an accident.
\item[Planck] (\texttt{spectrum\_is\_planck}). The photospheric blackbody
  seen at the orbit,
  %
  \begin{equation}
  F_\nu = \pi B_\nu(T_{\rm eff})\left(\frac{R_\star}{a}\right)^{2},
  \qquad
  B_\nu = \frac{2h\nu^{3}/c^{2}}{e^{h\nu/kT_{\rm eff}}-1},
  \label{eq:rad:planck}
  \end{equation}
  %
  converted to a flux per unit photon energy by $F_E = F_\nu\,{\rm d}\nu/{\rm d}E$
  with ${\rm d}\nu/{\rm d}E = 2.417989242\times 10^{14}$\,Hz\,eV$^{-1}$.
  For $h\nu/kT_{\rm eff} > 700$ the Wien form replaces $1/(e^{x}-1)$ so that
  no exponential overflows. The absolute scale comes from $T_{\rm eff}$ and
  $R_\star$, not from the two luminosities, so the grid-integrated flux
  generally differs from $J_{\rm XUV}$; \texttt{write\_setup\_report.f90}
  states both.
\item[Load] (\texttt{do\_read\_sed}). A measured spectral energy
  distribution read by \texttt{read\_sed} (\texttt{sed\_read.f90}): two
  columns, bin central wavelength in \AA{} and flux at the planet in
  erg\,cm$^{-2}$\,s$^{-1}$\,\AA$^{-1}$, ascending in wavelength. Between
  two rows the field is interpolated log-linearly in $\log E$, exact for the
  power-law segments a tabulated continuum is made of; outside the band it
  returns zero, and \texttt{spectrum\_covers\_eV} must be asked first so
  that the run stops rather than integrating a zero the input never stated.
\end{description}

A monochromatic option (\texttt{is\_monochr}) puts the whole
$10^{L_{\rm EUV}}/4\pi a^{2}$ at a single energy $E_{\rm low}$.

\paragraph{Bands of the photon grid.} \texttt{set\_energy\_vectors.f90}
lays the grid out over at most six bands,
\[
[E_{\rm floor},E_{\rm abs}],\;
[E_{\rm abs},13.598\,{\rm eV}],\;
[13.598,24.587],\;
[24.587,54.418],\;
[54.418,E_{\rm mid}],\;
[E_{\rm mid},E_{\rm top}],
\]
with fixed bin budgets $(20,89,50,50,50,50)$. $E_{\rm low}$, $E_{\rm mid}$
and $E_{\rm top}$ are stated by the
\texttt{[E\_low,E\_mid,E\_high]} line of \texttt{input.inp} (the tutorial
run uses $13.60$, $123.98$ and $1.24\times10^{3}$\,eV; with
\texttt{Use only EUV? True} the top is held at $1.24\times10^{3}$\,eV and
the X-ray normalization is zero). $E_{\rm abs}$ is
\texttt{sub\_lyman\_absorber\_threshold\_eV()}, the lowest ionization
threshold of an absorber of the bulk gas below the hydrogen edge: the
He\,$2^3$S metastable at 4.768\,eV when it is tracked, or the neutral stage
of an active low-IP metal (K\,{\sc i} 4.341, Na\,{\sc i} 5.139,
Ca\,{\sc i} 6.113, Mg\,{\sc i} 7.646, Fe\,{\sc i} 7.902, Si\,{\sc i} 8.152,
S\,{\sc i} 10.36\,eV) when one is lower. $E_{\rm floor}$ is
\texttt{photon\_grid\_floor\_eV()}, which lowers that further to
$E_{\rm th}({\rm H}, n{=}2) = E_{\rm th}({\rm H})/4 = 3.3996$\,eV, the head
of the Balmer continuum, whenever the excited-hydrogen coupling is armed.
The interval $[E_{\rm floor},E_{\rm abs}]$ carries no bulk-gas absorber: it
exists so that the H($n{=}2$) band is stated, that rate being a closed-form
integral over it (Sec.~\ref{sec:ionization:n2}) and not a sum over these bins.

\paragraph{Every ionization threshold is a bin edge.} A photoionization
rate is the rectangle sum $\sum_k F_k \sigma_k \Delta E_k/E_k$, so a bin
that \emph{contains} a threshold would charge the absorber its cross
section, read at the bin center, over the part of the bin below the threshold
where it cannot absorb at all. \texttt{ionization\_thresholds\_active}
therefore collects every threshold the run carries (H\,{\sc i},
He\,{\sc i}, He\,{\sc ii}, He\,$2^3$S, H$_2$ and its two channel
thresholds, and every photoionizable stage of an element with nonzero
abundance together with its subshell turn-ons); each band is cut at those
falling strictly inside it, and its budget is shared over the subintervals in
proportion to their logarithmic widths (largest remainder, at least one bin
each). The bins are geometric and $E_k$ is the \emph{geometric} mean of the
edges, the more accurate abscissa of a rectangle rule for a power-law
integrand. A grid whose points sat \emph{on} the thresholds instead overstates
the exact $P_{\rm HI}$, $P_{\rm HeI}$ and $P_{\rm HeII}$ by factors
$1.095$, $1.016$ and $1.030$.

\paragraph{Dayside dilution.} Every stellar beam of the code is multiplied by
the dilution $\xi_{\rm day}$ of Table~\ref{tab:2d}: the XUV grid in
\texttt{set\_energy\_vectors}, $F_{\rm XUV} \to \xi_{\rm day}F_{\rm XUV}$;
the stellar Ly$\alpha$ beam in \texttt{lya\_rt.f90} and
\texttt{excited\_hydrogen.f90}; and the four far-ultraviolet bands through
\texttt{fuv\_band\_flux}.

\paragraph{Bands stated by their integrated flux.} Four far-ultraviolet
bands lie almost entirely below the photon grid and enter as
band-integrated energy fluxes at the planet's orbit rather than through
the grid, in the band order of \texttt{water\_photolysis.f90}: LW
(912--1201\,\AA), B2 (the H\,{\sc i} Ly$\alpha$ line at 1215.67\,\AA,
band 1202--1230\,\AA), B3 (1231--1450\,\AA) and B4 (1451--2304\,\AA). The
edges are fixed by the physics of the absorbers and not chosen: 912\,\AA{}
is the Lyman edge, 1201\,\AA{} the red end of the H$_2$ Lyman and Werner
line list, and the B2/B3/B4 edges follow the H$_2$O branching ratios
(Sec.~\ref{sec:mol:oxygen}). Each band flux has one of two sources. A key
in \texttt{input.inp} states it (\texttt{Stellar LW flux},
\texttt{Stellar Lya flux}, \texttt{Stellar FUV B3 flux},
\texttt{Stellar FUV B4 flux}), and a stated value, zero included, always
stands. Where no key is given and the spectrum is a loaded file
(\texttt{Spectrum type: Load}), the three continuum bands are integrated
from that file by \texttt{sed\_band\_integrated\_flux} (\texttt{sed\_read.f90}),
\begin{equation}
F_{\rm band} = \int_{\lambda_{\rm lo}}^{\lambda_{\rm hi}} F_\lambda\,d\lambda,
\label{eq:rad:bandflux}
\end{equation}
the exact integral of the piecewise-linear table, the flux interpolated at
each clipped band edge before the trapezoid (a constant or a linear
spectrum integrates exactly whatever the nodes), the file being already at
the planet and read once for the three bands: LW on any molecular run, B3
and B4 when the oxygen chemistry is on. A file that does not reach both
edges of a band the run needs stops the run and asks for the key. The two analytic spectrum types cannot supply these
bands (a power law is the XUV fit extrapolated downward, a Planck curve
the photosphere, while the stellar far-ultraviolet is chromospheric), and
the ionizing grid retained by \texttt{read\_sed} ends at the lowest active
threshold, so the file is re-read for them. Band B2 is never integrated:
a Ly$\alpha$ line flux reconstructed from observations is a better number
than a trapezoid over whatever rows the file has across 30\,\AA, so it is
always a key. \texttt{fuv\_band\_flux} returns $\xi_{\rm day}F_{\rm band}$,
the dilution being applied where the beam is used so that the stated or
integrated values keep their meaning, the band flux at the planet's orbit;
the setup report states all four and names the source of each.

\subsection{Photoionization cross sections}
\label{sec:radiation:crosssec}

All cross sections are carried in units of $10^{-18}$\,cm$^{2}$
(\texttt{cross\_sec.f90}) and are evaluated once, at the grid centers, into
the vectors \texttt{s\_hi}, \texttt{s\_hei}, \texttt{s\_heii},
\texttt{s\_heiTR}, \texttt{s\_h2} and the table \texttt{sigma\_tab}.

\begin{description}
\item[Hydrogenic ions] (H\,{\sc i}, He\,{\sc ii}), function
  \texttt{sigma}. With $\epsilon = \sqrt{E/E_{\rm th}-1}$,
  \begin{equation}
  \sigma(E) = \frac{6.3}{Z^{2}}\left(\frac{E_{\rm th}}{E}\right)^{4}
  \frac{\exp\!\left[4-4\arctan(\epsilon)/\epsilon\right]}
       {1-\exp(-2\pi/\epsilon)},\qquad E>E_{\rm th},
  \label{eq:rad:hydrogenic}
  \end{equation}
  and $6.3/Z^{2}$ exactly at threshold, zero below it. $E_{\rm th}$ is the
  \emph{measured} potential of the ion, the same energy the grid uses as a
  band edge and the photoelectron is charged against.
\item[He\,{\sc i} ground singlet] The \citet{Verner1996} single-shell fit,
  \begin{equation}
  x = \frac{E}{E_0}-y_0,\quad z=\sqrt{x^{2}+y_1^{2}},\quad Q=5.5-P/2,
  \qquad
  \sigma = \sigma_0\left[(x-1)^{2}+y_w^{2}\right]z^{-Q}
           \left(1+\sqrt{z/y_a}\right)^{-P},
  \label{eq:rad:vfky96}
  \end{equation}
  with the published parameters $(E_0,\sigma_0,y_a,P,y_w,y_0,y_1) =
  (13.61, 949.2, 1.469, 3.188, 2.039, 0.4434, 2.136)$ and $E_{\rm th}$ taken
  as the global He\,{\sc i} threshold (it only gates the fit). The legacy
  ATES two-term fit $0.6935/[(E/100)^{1.82}+(E/100)^{3.23}]$ is restored by
  \texttt{ates\_photoion\_rate} (default off).
\item[He\,$2^3$S metastable] Two forms of \eqref{eq:rad:vfky96} joined by a
  log-linear bridge: wing A from the 4.768\,eV threshold to the Cooper
  minimum at 34.70\,eV, wing B from the resonance-averaged bump at 45.59\,eV
  upward, and a power law joining them continuously in between. Each wing is
  a fit at the few-percent level to the \citet{Norcross1971} and TOPbase
  background.
\item[Metal ions] \texttt{metal\_photoion\_sigma} returns the \emph{shell
  sum} over the seven subshells of \citet{Verner1996} and
  \citet{Verner1995}; the
  1995 form is \eqref{eq:rad:vfky96} with $y_0=y_1=0$ and
  $Q = 5.5+\ell-P/2$. Above the $E_{\max}$ of the 1996 outer-shell table the
  handover to the 1995 subshell fit is made, as in Verner's own
  \texttt{phfit2}; extrapolating the outer-shell fit is not conservative
  (the Fe\,{\sc i} row rises as $E^{2-Q}$ and reaches 24.9\,Mb at 1240\,eV
  against a true 0.51\,Mb). An inner-shell absorption is charged to one
  ionization of the absorbing stage with a photoelectron of
  $h\nu - E_{\rm th}({\rm outer})$, which approximates the Auger cascade and
  undercounts the charge state by at least one stage; this is the
  approximation the published exoplanet photochemistry models make
  \citep{Cecchi2009, Locci2022}. Fluorescence is neglected: with $E_{\rm top}=1240$\,eV only
  C, N, O and Na\,{\sc i} can have a K hole, and their K yields \citep{Krause1979} are 0.28, 0.52, 0.83 and 2.3
  percent.
\item[H$_2$] \texttt{sigma\_H2} is a piecewise \emph{measured} total
  photoabsorption: $k_{\rm bx}\times$ \citet{Backx1976} Table 1 over
  15.4--18\,eV with $k_{\rm bx}=1.016415$, \citet{Samson1994} Table 1
  over 18--300\,eV, and $k_{\rm hi}\times$ the \citet{Yan1998} eq.~(19) $E^{-7/2}$ sum-rule tail above 300\,eV with
  $k_{\rm hi}=0.885805$, all interpolated in $\log\sigma$ against $\log E$.
  No piece of the Yan fit survives below 300\,eV: its three published pieces
  do not meet at their own join energies, stepping by a factor 1.719 at
  18\,eV and 0.619 at 85\,eV.
\item[Opacity models] \texttt{opacity\_models.f90} dispatches the cross
  sections through a model selected in the optional \texttt{opacity.inp}:
  \texttt{A} analytic (the curves above, the default), \texttt{C} constant
  above threshold, \texttt{P} constant with the \citet{Robinson2012}
  pressure broadening $f(p) = 1 + a\,(p/p_{\rm piv})^{n}$, and \texttt{T} a
  two-column table. In model \texttt{P} the multiplier $f(p)$ of a cell is
  stored in \texttt{opa\_pf}$(j)$ and carried identically in the optical
  depth and in the local absorption, so that the photons a cell removes from
  the beam are the photons it absorbs.
\end{description}

\subsection{Transfer along the radial ray}
\label{sec:radiation:transfer}

\paragraph{The continuous form.} What the code evaluates is the
frequency-resolved absorption of the incident beam along the substellar
radial ray. With $F_E(a)$ the incident flux per unit photon energy at the
orbital distance $a$ (for the power law, $J_{\rm XUV} = L_{\rm XUV}/4\pi a^{2}$
of \eqref{eq:rad:bandsplit}; for the Planck and table types,
Sec.~\ref{sec:radiation:spectrum}), the flux reaching radius $r$ is
\begin{equation}
F_E(r) = F_E(a)\,e^{-\tau_E(r)},
\qquad
\tau_E(r) = \sum_a \sigma_{a,E}\int_r^{r_{\rm top}} n_a(r')\,dr',
\label{eq:rad:continuous_tau}
\end{equation}
the sum running over every photoionizing absorber $a$ and the column
starting at zero at the outer edge of the grid, $r_{\rm top} = r(N+N_g)$:
no absorber is carried beyond the domain, which is where these expressions
differ from Eq.~(4) of \citet{Caldiroli2021}, whose integral runs to $a$.
The photoionization rate of one particle of absorber $a$ and the
photoheating rate of the gas are then
\begin{equation}
\Gamma_a(r) = \int_{E_{{\rm th},a}}^{\infty}
   \frac{F_E(r)}{E}\,\sigma_{a,E}\,dE,
\qquad
\mathcal{H}(r) = \sum_a n_a(r)\int_{E_{{\rm th},a}}^{\infty}
   F_E(r)\left(1-\frac{E_{{\rm th},a}}{E}\right)\sigma_{a,E}\,dE,
\label{eq:rad:continuous_rates}
\end{equation}
which are Eq.~(3) of \citet{Caldiroli2021} and its rate counterpart,
$1-E_{\rm th}/E$ being the fraction of the photon energy carried off by the
photoelectron. Two factors the code adds to \eqref{eq:rad:continuous_tau}
are stated where they enter below: the pressure-dependent opacity weight
$f(p)$ of the optional model~\texttt{P}, and the secondary-ionization heat
fraction $f_{\rm heat}(E)$ that multiplies the photoelectron share in
$\mathcal{H}$ when secondary ionization is on
(Sec.~\ref{sec:radiation:secondary}). The rest of this section is the
discretization of \eqref{eq:rad:continuous_tau} and
\eqref{eq:rad:continuous_rates} on the photon grid and on the radial cells,
each line of which was read against \texttt{photoionization\_field\_at\_cell\_HHe}
(\texttt{util\_ion\_eq.f90}), the one routine that forms the field and the
rates of a cell.

The transfer is frequency-integrated over the grid of
Sec.~\ref{sec:radiation:spectrum} and is pure absorption: there is no
scattering term and no diffuse field in the XUV. For cell $j$ the column of
each absorber \emph{outside} the cell is accumulated from the top down by
\texttt{advance\_starward\_columns} (\texttt{util\_ion\_eq.f90}) on the same
rectangle rule the grid uses,
\begin{equation}
N_a^{\rm out}(j) = \sum_{j'>j} n_a(j')\,\Delta r_{j'}R_0\,f(p_{j'}),
\label{eq:rad:column}
\end{equation}
and the cell's \emph{own} depth is formed from its own density rather than
as a difference of two columns, which is the same number in exact
arithmetic and loses its digits to cancellation for a thin cell. The two
optical depths at photon energy $E_k$ are
\begin{equation}
\tau^{\rm out}_k = \sum_a \sigma_{a,k}N^{\rm out}_a,
\qquad
\Delta\tau_k = \sum_a \sigma_{a,k}\,n_a(j)\,\Delta r_j R_0 f(p_j),
\label{eq:rad:tau}
\end{equation}
the sum running over H\,{\sc i}, the He\,{\sc i} ground singlet,
He\,{\sc ii}, He\,$2^3$S when tracked, H$_2$ when the molecular chemistry is
on, and every photoionizable metal stage.

The attenuated flux the cell's \emph{rate} sees is the \emph{mean} of the
beam over the cell, not its value at either face
(\texttt{cell\_mean\_attenuation} in \texttt{utilities.f90}). Since the
densities are uniform inside a cell, $\tau(s) = \tau^{\rm out}+s\,\Delta\tau$
and
\begin{equation}
\left\langle e^{-\tau}\right\rangle_{\rm cell}
 = e^{-\tau^{\rm out}}\,\frac{1-e^{-\Delta\tau}}{\Delta\tau},
\label{eq:rad:cellmean}
\end{equation}
the second factor being \texttt{absorbed\_fraction\_per\_unit\_depth}, which
switches to its series $1-\Delta\tau/2+\Delta\tau^{2}/6$ below
$\Delta\tau = 10^{-8}$. Only this mean makes the photons the cell absorbs
equal the photons the beam loses across it; a one-point rule at the inner
face is low by 42 percent at $\Delta\tau=1$, the depth a cell carries at the
ionization front. The optional \texttt{alpha} form of the 2-D approximation replaces
\eqref{eq:rad:cellmean} by \eqref{eq:alpha2d}, which has no closed form.

The working flux of the cell is therefore
\begin{equation}
\mathcal{F}_k(j) = F_{{\rm XUV},k}\,
   \left\langle e^{-\tau}\right\rangle_{k,{\rm cell}}\, f(p_j),
\label{eq:rad:workingflux}
\end{equation}
the model-\texttt{P} factor being folded once into the shared flux weight
rather than into each cross section, which is identical because every
integrand below is linear in exactly one cross section.

\subsection{Photoionization and photoheating rates}
\label{sec:radiation:rates}

Every integral is the rectangle sum over the grid,
$\sum_k(\cdot)\,\Delta E_k$, with $\Delta E_k$ the exact bin width, so
$\sum_k \Delta E_k = E_{\rm top}-E_{\rm floor}$ exactly. For absorber $a$
the photoionization rate per particle and the photoheating rate of
\emph{one} particle of that absorber are
\begin{align}
P_a &= \sum_k \mathcal{F}_k\,\frac{\sigma_{a,k}}{E_k}\,\Delta E_k
       \times 10^{-18}\,\mathrm{erg^{-1}eV},
  \label{eq:rad:photorate}\\
h^{(1)}_a &= \sum_k \mathcal{F}_k\,
   \left(1-\frac{E_{{\rm th},a}}{E_k}\right) f_{\rm heat}(E_k)\,
   \sigma_{a,k}\,\Delta E_k \times 10^{-18}\ \ [\mathrm{erg\,s^{-1}}],
  \label{eq:rad:photoheat}
\end{align}
where the factor $1-E_{\rm th}/E$ is the share of the photon left as
photoelectron kinetic energy (\texttt{photoelectron\_share}; the key
\texttt{photoheat\_photon\_fraction}, default off, replaces it by a constant
fraction), and $f_{\rm heat}$ is the secondary-ionization heat fraction of
Sec.~\ref{sec:radiation:secondary}. The heating of the cell is the contraction
of these one-particle rates with the composition,
\begin{equation}
\mathcal{H} = \sum_a n_a\,h^{(1)}_a,
\label{eq:rad:heatcontraction}
\end{equation}
performed in exactly one place (\texttt{photoheating\_of\_cell}) and split
by absorber into the six columns H\,{\sc i}, He\,{\sc i}, He\,{\sc ii},
He\,$2^3$S, H$_2$, metals. The heating efficiency written to the output is
\begin{equation}
\eta_{\rm heat} = \frac{\mathcal{H}}{\displaystyle\sum_k \mathcal{F}_k
        \Big(\sum_a \sigma_{a,k} n_a\Big)\Delta E_k\times 10^{-18}},
\label{eq:rad:eta}
\end{equation}
the ratio of the energy deposited to the energy absorbed, set to zero in a
transparent or unilluminated cell.

\paragraph{Molecular channels.} H$_2$ absorption is resolved into four
mutually exclusive final states whose cross sections are \emph{shares} of
$\sigma_{{\rm H}_2}$, so the opacity and the total H$_2$ destruction rate are
unchanged. Each channel is charged its own threshold in
\eqref{eq:rad:photoheat}, and the two deposits that are not a photoelectron,
the kinetic energy release of the double channel and the internal excitation
the neutral channel radiates away, are stated with the channels themselves in
Sec.~\ref{sec:mol:h2photo}.

\paragraph{Self-consistency of the field with the composition.} Nothing in
\eqref{eq:rad:tau} depends on any cell inside $j$, which is what lets the
equilibrium sweep march outside in and re-form the field of one cell against
its own composition (Sec.~\ref{sec:ionization:solve}); the cell mean
\eqref{eq:rad:cellmean} is formed once per pass so that no second attenuation
is stacked on the first. The rates are constants of the cell's nonlinear
solve, carrying the field dependence inside the residual meaning a
re-integration of the photon grid at every residual evaluation.

\subsection{Secondary ionization by photoelectrons}
\label{sec:radiation:secondary}

A photoelectron faster than $E_{\rm sec} = 30$\,eV does not thermalize
completely: part of its energy ionizes further neutrals.
\texttt{electron\_energy\_degradation.f90} divides the labor between two
published determinations. \citet{Shull1985} supply
the energy budget of the ionization channels as a function of the ionized
fraction $x$,
\begin{equation}
\phi_{\rm HI}(x) = 0.3908\left(1-x^{0.4092}\right)^{1.7592},
\qquad
\phi_{\rm HeI}(x) = 0.0554\left(1-x^{0.4614}\right)^{1.6660},
\label{eq:rad:svs85}
\end{equation}
whose form vanishes at $x=1$ where no neutral is left. \citet{Dalgarno1999}
supply the heating fraction and everything that
depends on molecular hydrogen; their Table 4 mean energies per ion pair,
$W(x,E_0) = a(1+c\,x^{b})$, give the energy dependence,
\begin{equation}
f_{\rm HI}(x,E_0) = \phi_{\rm HI}(x)\frac{W_{\rm H^+}(x,1\,{\rm keV})}
                                          {W_{\rm H^+}(x,E_0)},
\qquad
f_{\rm HeI}(x,E_0) = \phi_{\rm HeI}(x)\frac{W_{\rm He^+}(x,1\,{\rm keV})}
                                            {W_{\rm He^+}(x,E_0)},
\label{eq:rad:sec_energy}
\end{equation}
so that the high-energy limit is exactly \citeauthor{Shull1985}'s. The
division of that energy among the targets is recomputed on the cell's own
densities, using the collision weight $1.89$ that \citeauthor{Dalgarno1999}'s
eqs.~(9) and (10) imply (their coefficients $1.89$ and $0.53$ are reciprocal
to 0.17 percent), with $r_{\rm ref} = n({\rm He\,I})/n({\rm H\,I}) = 0.1$
the reference composition of the fits:
\begin{equation}
w = (n_{\rm HI}+1.89\,n_{\rm H_2})r_{\rm ref}f_{\rm HI} + n_{\rm HeI}f_{\rm HeI},
\quad
c_{\rm HI} = \frac{f_{\rm ion}\,r_{\rm ref}f_{\rm HI}}{w},
\quad
c_{\rm H_2} = 1.89\,c_{\rm HI},
\quad
c_{\rm HeI} = \frac{f_{\rm ion}f_{\rm HeI}}{w},
\label{eq:rad:sec_split}
\end{equation}
with $f_{\rm ion} = f_{\rm HI}+f_{\rm HeI}$. These coefficients are
\emph{per target particle} (units cm$^{3}$), so
$n_{\rm HI}c_{\rm HI}+n_{\rm H_2}c_{\rm H_2}+n_{\rm HeI}c_{\rm HeI}=f_{\rm ion}$
identically and nothing is divided by a vanishing neutral density. The heat
fraction is
\begin{equation}
f_{\rm heat}(E_0) = 1 - \left[1-\eta_{\rm mix}(x,E_0)\right](1-x)
                     + f_{\rm diss} + f_{\rm vib} + f_{\rm fluor},
\label{eq:rad:fheat}
\end{equation}
with $\eta_{\rm D} = 1 + (\eta_{\rm D,0}-1)/(1+c\,x^{b})$ their eq.~(14) and
$\eta_{\rm mix}$ their eq.~(12) mixing of the H/He and H$_2$/He efficiencies.
The factor $(1-x)$ is this code's closure and not either source's: the
channels that compete with heat need a neutral to act on, and
\citeauthor{Dalgarno1999}'s form leaves 10 percent of the primary energy unthermalized even at
$x=1$. The three molecular terms (fragment kinetic energy, direct
vibrational excitation of $v=1,2$, and the vibrational energy the B and C
states return on fluorescing) vanish identically with $n({\rm H}_2)$; the
last two carry the collisional quench fraction of
\texttt{h2\_vibrational\_relaxation.f90}, since a radiated excitation is not
heat.

The secondary rates are added to the primary ones of
\eqref{eq:rad:photorate}: for an absorber $a$ whose photoelectron energy
exceeds $E_{\rm sec}$,
\begin{equation}
P^{\rm sec}_{\rm HI} = \sum_k \mathcal{F}_k\,
   \frac{\sigma_{a,k}n_a}{E_k}\,c_{\rm HI}(E_k)\,
   \frac{E_k-E_{{\rm th},a}}{E_{\rm th}({\rm H\,I})}\,\Delta E_k
   \times 10^{-18}\,\mathrm{erg^{-1}eV},
\label{eq:rad:secrate}
\end{equation}
and likewise for He\,{\sc i} and H$_2$; the H$_2$ secondaries additionally
make one proton per 22 H$_2^+$ ions \citep[after their eq.~10]{Dalgarno1999}, which is an added yield and not a share.

\paragraph{Staged activation.} The key \texttt{Secondary\_ionization} is
\emph{on} by default (\texttt{use\_sec\_ion = .true.}) but \emph{staged}:
the runtime flag \texttt{sec\_ion\_active} stays false until the wind has
first converged without the coupling, because from a cold initial condition
the base feedback amplifies the startup transient into a runaway
(\texttt{EXHALE\_main.f90}, the block that sets \texttt{sec\_ion\_active =
.true.} once any of the mass-flux, step-change, plateau or steady gates is
met, and again before the JFNK handover so that Newton is never asked to
solve a system the marched state does not approximately satisfy).
\texttt{Secondary\_ionization: Immediate} applies it from step 0 and
\texttt{False} disables it entirely. Photoelectrons between 13.6 and 30\,eV
are still treated as depositing all their energy as heat, which is a limit
of the published tables and not a physical threshold.

\subsection{Reprocessed and non-stellar fields}
\label{sec:radiation:internal}

Three fields in the code are not the direct stellar beam.

\paragraph{Recombination radiation.} The ionizing photons emitted when
H\,{\sc ii}, He\,{\sc ii} and He\,{\sc iii} recombine are absorbed on the
spot (\texttt{recombination\_radiation\_absorbed} in
\texttt{util\_ion\_eq.f90}): each cell keeps the fraction $1-e^{-\tau}$ of
every channel at its own optical depth and shares it among H\,{\sc i},
He\,{\sc i}, He\,{\sc ii}, H$_2$ and the metals in the ratio of their
absorption coefficients; the part re-absorbed by the recombined species
itself is the on-the-spot cancellation of the net coefficients
\eqref{eq:ion:netalpha}. The H\,{\sc ii} ground captures
(\texttt{H\_rec\_escape}, default on) can be taken only by H\,{\sc i} and
the low-potential metals. With \texttt{He\_rec\_coupling} (default on)
each He\,{\sc ii} cascade exit is represented at one photon energy: the
ground-capture continuum starting at $E_{\rm th}({\rm He\,I})$, the
$2^1P\to1^1S$ resonance line at 21.2\,eV, the $2^3S\to1^1S$ line at
19.8\,eV, and the part of the $2^1S$ two-photon continuum above the hydrogen
edge. The last is not a line: integrating the \citet{Drake1969} shape over
13.598--20.62\,eV gives 0.5564 ionizing photons per decay carrying
8.9646\,eV, a mean in-band photon energy of 16.110\,eV and hence 2.512\,eV of
photoelectron energy per photon, against 3.511\,eV for a flat distribution
over the same window. The He\,{\sc iii} channels are the ground capture at
54.4\,eV, the direct capture into He\,{\sc ii} $n=2$ \citep{Seaton1959},
whose continuum starts at 13.605\,eV, He\,{\sc ii} Ly$\alpha$ at 40.81\,eV,
and the He\,{\sc ii} $2s$ two-photon continuum, whose spectrum is the fit of
\citet{Nussbaumer1984} carried in three bands (13.6--15.4, 15.4--24.6 and
above 24.6\,eV: 0.112, 0.577 and 0.736 photons per decay). Every case-B
capture of He\,{\sc iii} ends in $2s$ (the fraction of
\citealt{Pengelly1964}) or $2p$, and $2s$ decays by two photons at
526.823\,s$^{-1}$ (\citealt{Drake1986} eq.~27 at $Z=2$, with the
$\alpha$-particle mass) unless a proton or He$^{2+}$ collision moves it to
$2p$ first \citep{Pengelly1964b}.

\paragraph{Lyman-Werner band.} The 912--1201\,\AA{} beam is solved once for
all of its absorbers together (\texttt{fuv\_lw\_photon\_field}), because an
H$_2$O molecule cannot absorb a photon an H$_2$ line has already taken. Its
transmission down to a face is $T = (1-A)e^{-\tau_c}$ with
$A = \int \sigma_{\rm pump}f_{\rm shield}(N_{\rm H_2})n_{\rm H_2}{\rm d}r$
summed along the path, each cell contributing at its own $(T, n)$ the same
integral its own line rate takes (the table is built for a uniform slab, and
reading $A$ of the whole column at the temperature of the cell the beam has
reached would count as taken photons the lines of the cells above have not
taken, or the reverse),
and across one cell the identity
\begin{equation}
T^{\rm out}-T^{\rm in} = e^{-\tau^{\rm out}}
  \left[(1-A^{\rm out})\left(1-e^{-\Delta\tau}\right)
        + \Delta A\, e^{-\Delta\tau}\right]
\label{eq:rad:lwsplit}
\end{equation}
is exact, so the photons the cell removes split into a continuum share and a
line share with nothing left over. The continuum factor $e^{-\tau_c}$ is the
\citet{Draine1996} eq.~(40) term, identically 1 where the oxygen option is
off. Line self-shielding is the tabulated function of
$(T, n, N_{\rm H_2})$ of Sec.~\ref{sec:mol:lw}. CO is shielded on the same beam by
the \citet{Visser2009} function
$\Theta(N_{\rm CO},N_{\rm H_2})$ (\texttt{co\_self\_shielding\_table.f90}).

\paragraph{Ly$\alpha$ and the Balmer continuum} are described with the
excited-hydrogen model in Sec.~\ref{sec:ionization:n2}.
